% !TeX spellcheck = ru_RU
% !TEX root = vkr.tex

\section*{Заключение}

Таким образом, в рамках преддипломной практики были достигнуты следующие результаты:
\begin{itemize}
    \item разработан прототип анализатора, который реализует следующие анализы;
          \begin{itemize}
              \item поиск входных и выходных операндов инструкций;
              \item определение, обращается ли инструкция к памяти;
          \end{itemize}
    \item при помощи анализатора сравнены следующие спецификации RISC-V ISA:
          \begin{itemize}
              \item LLVM v19 против sail-riscv\footnote{\url{https://github.com/riscv/sail-riscv/commits/05b845c91d1c1db7b361fc8d06e815b54ca0b07a} (дата обращения: \DTMdate{2025-05-17}).};
              \item <<P>> extension (Version 0.9.11-draft-20211209\footnote{\url{https://tools.cloudbear.ru/docs/riscv-p-extension-0.9.11-20211209.pdf} (дата обращения: \DTMdate{2025-05-17}).}) в LLVM против sail-riscv-p\footnote{\url{https://github.com/umcann123/sail-riscv-p} (дата обращения: \DTMdate{2025-05-17}).};
                    \begin{itemize}
                        \item для sail-riscv-p дополнительно определены для каждой инструкции какие CSR регистры читает и пишет;
                    \end{itemize}
          \end{itemize}

\end{itemize}

Исходный код анализатора: \url{https://github.com/s-khechnev/llvm_td_lint_playground}.

Имя коммитера: s-khechnev.

% \textbf{Обязательно.}
% Список результатов, который будет либо один к одному соответствовать задачам из раздела~\ref{sec:task}, либо их уточнять (например, если было \enquote{выбрать}, то тут \enquote{выбрано то-то}).

% \begin{itemize}
%     \item Результат к задаче №1.
%     \item Результат к задаче №2.
%     \item и т.д.
% \end{itemize}
% \noindent Если работа на несколько семестров, отчитывайтесь только за текущий.
% Можно в свободной форме обрисовать планы продолжения работы, но не увлекайтесь~--- если работа будет продолжена, по ней будет ещё один отчёт.

% В заключении \emph{обязательна} ссылка на исходный код, если он выносится на защиту, либо явно напишите тут, что код закрыт.
% Если работа чисто теоретическая и это понятно из решённых задач, про код можно не писать.
% Обратите внимание, что ссылка на код должна быть именно в заключении, а не посреди раздела с реализацией, где её никто не найдёт.

% Старайтесь оформить программные результаты работы так, чтобы это был один репозиторий или один пуллреквест, правильно оформленный~--- комиссии тяжело будет собирать Ваши коммиты по всей истории.
% А если над проектом работало несколько человек и всё успело изрядно перемешаться, неизбежны вопросы о Вашем вкладе.

% Заключение люди реально читают (ещё до \enquote{основных} разделов работы, чтобы понять, что же получилось и стоит ли вообще работу читать), так что оно должно быть вылизано до блеска.
